\begin{lemma}[Proved: a length invariant of an exact cycle of acyclic complexes]
\label{C:three_cycle_length}
Let $X,Y,Z$ be acyclic two-periodic complexes in an abelian category
of finite-length objects. Suppose
\[
 X\longrightarrow Y\longrightarrow Z\longrightarrow X[1]
\]
is exact in each degree, with the first two maps of degree zero
and the last of degree one. Put
$h(C)=\ell\operatorname{im}(d_C^0)-\ell\operatorname{im}(d_C^1)$.
Then $h(Y)\equiv h(X)+h(Z)\pmod3$.
\end{lemma}
\begin{proof}
For any short exact sequence of two-periodic complexes
$0\to U\to V\to W\to0$, let $a_i$ be the length of the image
of its cohomological connecting map
$H^i(W)\to H^{i+1}(U)$. Lifting cycles and taking kernel lengths
gives
\[
 \ell\operatorname{im}d_V^i
  =\ell\operatorname{im}d_U^i+
    \ell\operatorname{im}d_W^i+a_i.
\]
Consequently $h(V)=h(U)+h(W)+a_0-a_1$.
If $V$ is acyclic, its connecting maps are isomorphisms and
$a_0-a_1=\ell H^1(U)-\ell H^0(U)=-\chi(U)$, where
$\chi(U)=\ell U^0-\ell U^1$.

Let $I,J,K$ be the images of the three maps in the displayed
cycle, so that $K\subset X[1]$. Exactness gives short exact
sequences of complexes
\[
 0\to K[-1]\to X\to I\to0,\quad
 0\to I\to Y\to J\to0,\quad
 0\to J\to Z\to K\to0.
\]
Their acyclic middle terms imply
$\chi(J)=-\chi(I)$ and $\chi(K)=\chi(I)$.
Applying the preceding formula, and using
$h(K[-1])=-h(K)$, yields
\[
 h(Y)-h(X)-h(Z)=-3\chi(I).
\]
This proves the congruence without asserting integer additivity.
\end{proof}

\begin{proposition}[Proved: a cubic-thickening invariant with values modulo three]
\label{C:cubic_Bockstein}
Let $A$ be a characteristic-zero DVR with maximal ideal $(2)$,
let $F=A/2A$, and put $k=A[\epsilon]/\epsilon^3$.
For a free $k$-representation $M$ of $C_2$ let
$V=\widehat H^0(C_2,M/\epsilon M)$ and
$W=\widehat H^1(C_2,M/\epsilon M)$.
The first $\epsilon$-Bockstein gives a natural two-periodic
differential $b$ on $V\oplus W$.
There is also a natural even operator $c$ commuting with $b$:
if $T(\epsilon)$ is any lift of the involution matrix to
$A[[\epsilon]]$, its action is induced by
\[
 C_0=\left.\frac{T(\epsilon)^2-1}{\epsilon^3}
                         \right|_{\epsilon=0}.         \tag{C.1.21}
\]
For every monic irreducible $q\in F[c]$ other than $c$, the
$q$-primary complex $(V_q,W_q,b)$ is acyclic, and
\[
 \nu_q(M)=\ell_{F[c]}\operatorname{im}(b^0:V_q\to W_q)
           -\ell_{F[c]}\operatorname{im}(b^1:W_q\to V_q)
                 \pmod3                               \tag{C.1.22}
\]
is additive on admissible short exact sequences.
\end{proposition}
\begin{proof}
Write $T=T_0+\epsilon T_1+\epsilon^2T_2$ modulo $\epsilon^3$.
The two-periodic Tate differential has constant terms $T_0-1$
and $T_0+1$, and the coefficient of $\epsilon$ is $T_1$.
Thus the Bockstein is induced by $T_1$, as in
Lemma~\ref{C:epsilon_Bockstein}. The coefficient identities are
\[
 T_0^2=1,\quad T_0T_1+T_1T_0=0,\quad
 T_0T_2+T_2T_0+T_1^2=0.
\]
The last identity makes $b^2=0$ on Tate cohomology: on a plus
cycle $T_1^2$ is minus a $(T_0+1)$-boundary, and on a minus
cycle it is minus a $(T_0-1)$-boundary.

We verify all assertions about $c$, including its independence
of choices. Put $C=(T^2-1)/\epsilon^3$ in the formal matrix ring.
The identity $TC=CT$ implies $T_0C_0=C_0T_0$, so $C_0$ acts on
the two Tate groups. Changing the lift $T$ by $\epsilon^3U$
changes $C_0$ by $T_0U_0+U_0T_0$. On a plus Tate cycle this
is a norm boundary, and on a minus cycle it is a difference
boundary; hence the induced operator is unchanged.

For a morphism of representations choose formal matrix lifts
$f,T_M,T_N$ and write
$T_Nf-fT_M=\epsilon^3Q$. Then
\[
 C_N f-f C_M=T_NQ+QT_M.
\]
At $\epsilon=0$ the right side is the appropriate Tate boundary
on either kind of cycle. This proves naturality, and also covers
changes of the chosen $k$-basis. The coefficient of $\epsilon$
in $TC=CT$ says
$[T_1,C_0]=-[T_0,C_1]$; on Tate cohomology its right side is
a boundary of the opposite parity. Therefore $b$ commutes with $c$.

For an admissible sequence choose an underlying $k$-linear
splitting and use upper block-triangular lifts
$T_E=\left(\begin{smallmatrix}T_A&X\\0&T_B\end{smallmatrix}\right)$.
The curvature is likewise upper block triangular, say with
off-diagonal block $Y$. The off-diagonal part of $T_EC_E=C_ET_E$
gives
\[
 C_A X-XC_B=T_A Y-YT_B.
\]
At $\epsilon=0$ this proves that $c$ commutes with the Tate
connecting map: a plus input has a difference boundary at the
minus target, and a minus input has a norm boundary at the plus
target. The Bockstein $b$ anticommutes with that connecting map
by the explicit calculation in Lemma~\ref{C:epsilon_Bockstein}.
Thus the exact Tate six-cycle is a cycle of two-periodic complexes
with commuting $F[c]$-actions. Taking primary parts preserves
its exactness.

It remains to prove acyclicity on the nonzero primary parts.
The decomposition (C.1.14) of $M/\epsilon M$ allows removal of
its regular summands. The basis lift and equivariant retraction
in Theorem~\ref{C:dualnumber_C2_ring} work unchanged with
$\epsilon^3=0$; their only requirement is nilpotence of
$\epsilon$. Such summands contribute no Tate cohomology.
For the complement choose $T_0=\operatorname{diag}(I,-I)$ and write
\[
 T_1=\begin{pmatrix}0&B\\D&0\end{pmatrix},\qquad
 T_2=\begin{pmatrix}P&E\\G&Q\end{pmatrix}.
\]
The coefficient equations give $P=-BD/2$, $Q=DB/2$;
these expressions are integral. In particular $\bar B\bar D$
and $\bar D\bar B$ vanish. Direct multiplication shows that
the off-diagonal blocks of $T_1T_2+T_2T_1$ cancel, and the
diagonal blocks are $BG+ED$ and $DE+GB$.
On the Tate groups put
$b_2=\left(\begin{smallmatrix}0&\bar E\\\bar G&0\end{smallmatrix}\right)$.
Then
\[
 b^2=0,\qquad bb_2+b_2b=c.
\]
On a $q$-primary part with $q\ne c$, the operator $c$ is
invertible. If $b_2$ does not preserve that part, first project
it there by the primary projector, which commutes with $b$.
Composing the resulting map with $c^{-1}$ gives a contracting
homotopy. Hence the primary complex is acyclic. Apply
Lemma~\ref{C:three_cycle_length} in the finite-length category
of $q$-primary $F[c]$-modules to the exact Tate cycle. It gives
exactly the additivity of (C.1.22).
\end{proof}

\begin{theorem}[Proved: odd torsion in the dyadic cubic-thickening kernel]
\label{C:cubic_three_torsion}
For $k=A[\epsilon]/\epsilon^3$ as above and every monic
irreducible $q\in F[u]$, $q\ne u$, lift its companion matrix
to an invertible $B\in M_d(A)$, where $d=\deg q$. Put
\[
 T_q=\begin{pmatrix}I_d&\epsilon I_d\\
                     \epsilon^2B&-I_d\end{pmatrix},
 \qquad\eta_q=[T_q]-d(1+\mathrm{sgn}).                  \tag{C.1.23}
\]
Then every $\eta_q$ has order exactly three, all of them vanish
generically and at every residue-field point, and they generate
a subgroup
\[
 \bigoplus_{q\ne u}\mathbb Z/3
 \ \subset\ \ker\bigl(R_k(C_2)\longrightarrow R_{k[1/2]}(C_2)\bigr).
                                                               \tag{C.1.24}
\]
Thus $\vartheta_k$ fails in degree two. This failure occurs even
though $3$ is a unit in $k$ and the acting group has order two.
This is a subgroup calculation, not a classification of the
entire cubic-thickening representation group.
\end{theorem}
\begin{proof}
The formal square of $T_q$ is
$I+\epsilon^3\operatorname{diag}(B,B)$, so it is an involution
over $k$, with curvature $c=\operatorname{diag}(\bar B,\bar B)$.
Its plus and minus Tate groups are both the simple module
$F[u]/(q)$; the first Bockstein is the identity from the minus
group to the plus group and zero in the other direction.
Consequently
$\nu_{q'}(\eta_q)=-\delta_{q,q'}\in\mathbb Z/3$.
This proves nonvanishing and will prove independence once the
order bound is established.

Write $M=(k^{2d},T_q)$ and
\[
 X=\begin{pmatrix}\epsilon^2B&-I_d\\
                    -\epsilon B&-\epsilon^2B\end{pmatrix},
 \qquad T_E=\begin{pmatrix}T_q&X\\0&T_q\end{pmatrix}.
\]
One checks $T_qX+XT_q=0$, so $E$ is an involution and
$0\to M\to E\to M\to0$ is admissible.
Label its four coordinate blocks $e_1,e_2,e_3,e_4$, each of
rank $d$. The vectors $v=e_4$ and
\[
 T_Ev=-e_1-\epsilon^2B e_2+\epsilon e_3-e_4
\]
generate a regular submodule $k[C_2]^d$. Together with $e_2,e_3$
they are an underlying $k$-basis, so the quotient is free.
On the quotient $e_4=0$ and
$e_1=-\epsilon^2B e_2+\epsilon e_3$; hence its matrix in the
ordered blocks $(e_3,e_2)$ is
\[
 Q_- =\begin{pmatrix}I_d&\epsilon^2I_d\\
                      -\epsilon B&-I_d\end{pmatrix}.
\]
Put $P=I-\epsilon
\left(\begin{smallmatrix}0&0\\B&0\end{smallmatrix}\right)$.
Then
\[
 P^{-1}Q_-P=Q_+,
 \qquad Q_+=\begin{pmatrix}I_d&\epsilon^2I_d\\
                         \epsilon B&-I_d\end{pmatrix}.
\]
The invertible matrix
$K=\left(\begin{smallmatrix}0&I_d\\-B&0\end{smallmatrix}\right)$
obeys $K(-T_q)=Q_+K$. Thus this quotient is
$M\otimes\mathrm{sgn}$, and
\[
 2[M]=d[k[C_2]]+[M\otimes\mathrm{sgn}].
\]
The evaluation sequence (C.1.7) also gives
$[M]+[M\otimes\mathrm{sgn}]=2d[k[C_2]]$.
Combining the two identities gives $3[M]=3d(1+\mathrm{sgn})$,
so $3\eta_q=0$. Its nonzero value under $\nu_q$ proves exact
order three, and the full family of $\nu_{q'}$ proves the
direct-sum assertion.

After inverting $2$ the two eigenspaces are free of rank $d$
over the local truncated-polynomial ring, so $\eta_q$ vanishes.
At every residue-field point $\epsilon=0$, and the same equality
is immediate. The invariants vanish on the two constant lines,
which span the degree-two polynomial image, proving the last
assertion about $\vartheta_k$.
\end{proof}
